\documentclass[10pt]{article}
\usepackage[utf8]{inputenc}
\usepackage[vietnamese]{babel}
\usepackage{amsmath,amssymb,amsfonts}
\usepackage{graphicx}
\usepackage{xcolor}
\usepackage{colortbl}
\usepackage{tabularx}
\usepackage{fancyhdr}
\usepackage[paperwidth=210mm, paperheight=280mm, left=18mm, right=18mm, top=20mm, bottom=20mm, headheight=14pt]{geometry}
\usepackage{tcolorbox}
\usepackage{newtxtext,newtxmath}

% Định nghĩa hệ màu sắc chuẩn phong cách giáo trình James Stewart
\definecolor{stewartcyan}{RGB}{0, 130, 165}
\definecolor{stewartred}{RGB}{185, 30, 45}
\definecolor{tableheadercolor}{RGB}{222, 235, 247}
\definecolor{margintextcolor}{RGB}{60, 60, 60}

% Cấu hình Header và Footer trang 80 (trang chẵn)
\pagestyle{fancy}
\fancyhf{}
\fancyhead[L]{\textsf{\textbf{80}}\hspace{2.5em}\textsf{\textbf{\color{stewartcyan}CHƯƠNG 2}}\hspace{1.2em}\textsf{Giới hạn và đạo hàm}}
\fancyhead[R]{}
\renewcommand{\headrulewidth}{0pt}

\begin{document}

\noindent
\begin{minipage}[t]{0.28\textwidth}
    \vspace*{80pt}
    
    % Bảng tính hệ số góc các cát tuyến
    \centering
    {\footnotesize
    \setlength{\tabcolsep}{4.5pt}
    \renewcommand{\arraystretch}{1.2}
    \begin{tabular}{|c|c|}
    \hline
    \rowcolor{tableheadercolor}
    $R$ & $m_{PR}$ \\
    \hline
    $(0, 10)$ & $-82{,}425$ \\
    $(0{,}02;\, 8{,}187)$ & $-74{,}200$ \\
    $(0{,}06;\, 5{,}488)$ & $-60{,}750$ \\
    $(0{,}08;\, 4{,}493)$ & $-55{,}250$ \\
    $(0{,}1;\, 3{,}676)$ & $-50{,}450$ \\
    \hline
    \end{tabular}
    }

    \vspace{155pt}

    % Nhãn Hình 5 bên lề trái (dóng hàng với đáy đồ thị Hình 5)
    \raggedright
    \noindent\textbf{\textsf{HÌNH 5}}

    \vspace{16pt}

    % Ghi chú vật lý bên lề trái (dóng hàng với công thức ước lượng độ dốc)
    \noindent{\footnotesize\color{margintextcolor}
    Ý nghĩa vật lý của kết quả trong Ví dụ~2 là cường độ dòng điện chạy từ tụ điện sang laser sau $0{,}04$ giây là khoảng $-65$ ampe.
    \par}

\end{minipage}%
\hfill
\begin{minipage}[t]{0.69\textwidth}
    \setlength{\parindent}{1.5em}
    
    \noindent Từ các điểm $P(0{,}04;\, 6{,}703)$ và $R(0, 10)$ trên đồ thị, ta tìm được hệ số góc của cát tuyến $PR$ là
    \begin{equation*}
    m_{PR} = \frac{10 - 6{,}703}{0 - 0{,}04} = -82{,}425
    \end{equation*}

    Bảng ở bên trái cho thấy kết quả của các phép tính tương tự đối với hệ số góc của các cát tuyến khác. Từ bảng này, ta kỳ vọng hệ số góc của tiếp tuyến tại $t = 0{,}04$ sẽ nằm trong khoảng giữa $-74{,}20$ và $-60{,}75$. Trên thực tế, giá trị trung bình của các hệ số góc của hai cát tuyến gần nhất là
    \begin{equation*}
    \frac{1}{2}(-74{,}20 - 60{,}75) = -67{,}475
    \end{equation*}
    Vì vậy, bằng phương pháp này, ta ước lượng hệ số góc của tiếp tuyến vào khoảng $-67{,}5$.

    Một phương pháp khác là vẽ một đường xấp xỉ với tiếp tuyến tại $P$ và đo độ dài các cạnh của tam giác $ABC$, như trong Hình~5.

    \begin{center}
        \vspace{-0.3em}
        \includegraphics[width=0.72\linewidth]{images/p0115_fig1.png}
        \vspace{-0.3em}
    \end{center}

    \noindent
    Điều này cho ta một ước lượng về hệ số góc của tiếp tuyến là
    \begin{equation*}
    -\frac{|AB|}{|BC|} \approx -\frac{8{,}0 - 5{,}4}{0{,}06 - 0{,}02} = -65{,}0 \tag*{\color{stewartcyan}\blacksquare}
    \end{equation*}

    \vspace{1.2em}
    \noindent
    {\color{stewartcyan}\rule{1.1ex}{1.1ex}}\hspace{0.5em}\textbf{\textsf{Bài toán vận tốc}}\quad
    Nếu quan sát đồng hồ tốc độ của một chiếc ô tô khi lái xe trong phố, bạn sẽ thấy rằng tốc độ không duy trì như cũ trong thời gian dài; nghĩa là, vận tốc của ô tô không phải là một hằng số. Khi nhìn vào đồng hồ tốc độ, chúng ta giả định rằng chiếc xe có một vận tốc xác định tại mỗi thời điểm, nhưng vận tốc ``tức thời'' được định nghĩa như thế nào?

    Chúng ta hãy xem xét \emph{bài toán vận tốc}: Tìm vận tốc tức thời của một vật chuyển động theo quỹ đạo thẳng tại một thời điểm cụ thể nếu biết vị trí của vật tại mọi thời điểm. Trong ví dụ tiếp theo, chúng ta sẽ khảo sát vận tốc của một quả bóng đang rơi. Qua các thí nghiệm được tiến hành từ bốn thế kỷ trước, Galileo đã phát hiện ra rằng quãng đường rơi của một vật rơi tự do bất kỳ tỷ lệ thuận với bình phương thời gian rơi của nó. (Mô hình rơi tự do này bỏ qua sức cản của không khí.) Nếu quãng đường rơi sau $t$ giây được ký hiệu là $s(t)$ và đo bằng mét, thì (tại bề mặt Trái Đất) quan sát của Galileo được biểu diễn bởi phương trình
    \begin{equation*}
    s(t) = 4{,}9t^2
    \end{equation*}

\end{minipage}

\end{document}